\section{How to use}
This report template is designed to help you create a professional-looking report with ease. Below are the steps to get started:
\subsection{Lists}
You can create itemized lists as follows:
\begin{itemize}
  \item First item in an itemized list.
  \item Second item in an itemized list.
\end{itemize}
You can create enumerated lists as follows:
\begin{enumerate}
  \item First item in an enumerated list.
  \item Second item in an enumerated list.
\end{enumerate}
\subsection{Tables}
Note that you can use \texttt{\textbackslash FloatBarrier} to prevent floats from moving past a certain point in the document. In other words, all figures and tables before \texttt{\textbackslash FloatBarrier} will appear before it in the final output.
\begin{table}
  \centering
  \caption{Example of a simple, modern table}
  \label{tab:simple-modern}
  \begin{tabular}{lcc}
    \toprule
    Item & Quantity & Price (USD) \\
    \midrule
    Notebook & 2 & 12.50 \\
    Pen & 5 & 1.20 \\
    Backpack & 1 & 45.00 \\
    \midrule
    Total &  & 71.90 \\
    \bottomrule
  \end{tabular}
\end{table}

\begin{table}
\centering
\caption{Sample formulas using tblr}
\label{tab:formulas}
\begin{tblr}{
    colspec = {l X[c] X[c]},
    row{1} = {font=\bfseries},
    hlines, vlines
}
Formula & Expression & Notes \\
Pythagorean & $a^2 + b^2 = c^2$ & Right triangle relation \\
Quadratic roots & $x = \frac{-b \pm \sqrt{b^2 - 4ac}}{2a}$ & Discriminant $\Delta = b^2 - 4ac$ \\
Euler identity & $e^{i\pi} + 1 = 0$ & Fundamental constants \\
Derivative & $\frac{d}{dx} x^n = n x^{n-1}$ & $n \in \mathbb{N}$ \\
Integral & $\int_0^{\infty} e^{-x}\,dx = 1$ & Convergent improper integral \\
Summation & $\sum_{k=1}^n k = \frac{n(n+1)}{2}$ & Arithmetic series \\
\end{tblr}
\end{table}

\begin{table}
\centering
\caption{Comprehensive benchmark comparison}
\label{tab:complex-table}
\begin{tblr}{
    colspec = {l *{4}{X[c]}},
    rowhead = 2,
    hlines,
    vlines,
    cell{1}{1-5} = {bg=gray!20,font=\bfseries},
    cell{2}{1-5} = {bg=gray!10,font=\itshape},
    cell{3}{1} = {r=2}{font=\bfseries},
    cell{5}{1} = {r=3}{font=\bfseries},
    cell{8}{2} = {c=2}{bg=yellow!20},
}
Category & Impl A & Impl B & Impl C & Impl D \\
Metric & Time (ms) & Memory (MB) & Throughput (ops/s) & Energy (J) \\
Core Ops & 12.4 & 128 & 4200 & 0.84 \\
                 & 11.9 & 130 & 4500 & 0.81 \\
I/O Bound & 55.0 & 256 & 1100 & 2.10 \\
                 & 53.2 & 250 & 1300 & 2.05 \\
                 & 54.7 & 248 & 1250 & 2.00 \\
Summary & \SetCell{c=2}{Avg speed-up: 1.05x} & Peak: 4500 & Min energy: 0.81 \\
Notes & \SetCell{c=4,halign=l}{Time: lower is better; Throughput: higher is better; Energy: lower is better.} \\
\end{tblr}
\end{table}
\FloatBarrier
\begin{table}
    \centering
    \caption{Scaled wide performance matrix}
    \label{tab:wide-matrix}
    \resizebox{0.7\textwidth}{!}{%
    \begin{tabular}{lrrrrrrrrrrrr}
        \toprule
        Component & M1 & M2 & M3 & M4 & M5 & M6 & M7 & M8 & M9 & M10 & M11 & Avg \\
        \midrule
        Parser        & 12.3 & 11.9 & 12.1 & 12.0 & 11.8 & 12.4 & 12.2 & 11.7 & 12.5 & 12.3 & 12.0 & 12.11 \\
        Lexer         & 5.4  & 5.2  & 5.3  & 5.5  & 5.1  & 5.4  & 5.3  & 5.2  & 5.6  & 5.4  & 5.3  & 5.36  \\
        Optimizer     & 33.2 & 32.8 & 33.5 & 33.0 & 32.7 & 33.8 & 33.3 & 32.9 & 34.0 & 33.1 & 33.4 & 33.24 \\
        CodeGen       & 21.5 & 21.0 & 21.7 & 21.3 & 21.2 & 21.9 & 21.4 & 21.1 & 22.0 & 21.6 & 21.5 & 21.47 \\
        JIT           & 44.0 & 43.5 & 44.2 & 43.8 & 43.6 & 44.5 & 44.1 & 43.7 & 44.7 & 44.3 & 44.0 & 44.04 \\
        Profiler      & 8.9  & 8.7  & 8.8  & 8.6  & 8.5  & 9.0  & 8.8  & 8.6  & 9.1  & 8.9  & 8.8  & 8.80  \\
        GC            & 57.0 & 56.2 & 57.8 & 56.9 & 56.0 & 58.1 & 57.5 & 56.4 & 58.3 & 57.2 & 57.1 & 57.05 \\
        Scheduler     & 14.2 & 13.9 & 14.5 & 14.1 & 13.8 & 14.6 & 14.3 & 13.7 & 14.7 & 14.4 & 14.2 & 14.22 \\
        Net I/O       & 62.5 & 61.8 & 63.1 & 62.2 & 61.5 & 63.4 & 62.7 & 61.6 & 63.6 & 62.8 & 62.4 & 62.51 \\
        Disk I/O      & 71.0 & 70.2 & 71.9 & 71.1 & 70.0 & 72.3 & 71.5 & 70.4 & 72.6 & 71.7 & 71.2 & 71.35 \\
        Cache Layer   & 9.3  & 9.1  & 9.4  & 9.2  & 9.0  & 9.5  & 9.3  & 9.1  & 9.6  & 9.4  & 9.2  & 9.28  \\
        Auth Module   & 4.8  & 4.7  & 4.9  & 4.8  & 4.6  & 5.0  & 4.9  & 4.7  & 5.1  & 4.9  & 4.8  & 4.85  \\
        API Layer     & 18.6 & 18.2 & 18.9 & 18.4 & 18.1 & 19.1 & 18.7 & 18.3 & 19.3 & 18.8 & 18.5 & 18.64 \\
        Frontend      & 25.2 & 24.8 & 25.6 & 25.0 & 24.7 & 25.9 & 25.4 & 24.9 & 26.1 & 25.5 & 25.3 & 25.30 \\
        Monitoring    & 6.7  & 6.6  & 6.8  & 6.5  & 6.4  & 6.9  & 6.8  & 6.5  & 7.0  & 6.8  & 6.7  & 6.70  \\
        Alerting      & 3.4  & 3.3  & 3.5  & 3.4  & 3.2  & 3.6  & 3.5  & 3.3  & 3.7  & 3.5  & 3.4  & 3.45  \\
        Billing       & 11.2 & 10.9 & 11.4 & 11.1 & 10.8 & 11.6 & 11.3 & 10.9 & 11.7 & 11.5 & 11.2 & 11.24 \\
        Analytics     & 39.4 & 38.7 & 40.1 & 39.2 & 38.5 & 40.6 & 39.8 & 38.6 & 40.9 & 39.9 & 39.3 & 39.55 \\
        ML Inference  & 83.0 & 82.1 & 84.2 & 83.4 & 81.9 & 85.0 & 84.1 & 82.3 & 85.4 & 84.3 & 83.6 & 83.93 \\
        Security Scan & 52.1 & 51.4 & 52.9 & 52.0 & 51.0 & 53.3 & 52.5 & 51.3 & 53.6 & 52.7 & 52.2 & 52.36 \\
        \midrule
        Median        & 21.5 & 21.0 & 21.7 & 21.3 & 21.2 & 21.9 & 21.4 & 21.1 & 22.0 & 21.6 & 21.5 & 21.47 \\
        95th pct      & 82.0 & 81.2 & 83.1 & 82.4 & 80.9 & 84.0 & 83.2 & 81.3 & 84.4 & 83.3 & 82.6 & 82.60 \\
        Min           & 3.4  & 3.3  & 3.5  & 3.4  & 3.2  & 3.6  & 3.5  & 3.3  & 3.6  & 3.5  & 3.4  & 3.45  \\
        Max           & 83.0 & 82.1 & 84.2 & 83.4 & 81.9 & 85.0 & 84.1 & 82.3 & 85.4 & 84.3 & 83.6 & 84.52 \\
        \bottomrule
    \end{tabular}%
    }
    \vspace{10pt}

    \footnotesize Times shown in ms; lower is better. Columns M1-M11 represent sequential benchmark runs.
\end{table}
\subsection{Mathematic equations}
For inline math, you can wrap your equation in \texttt{\textbackslash( ...\textbackslash)}, like so: \(E=mc^2\). For displayed equations, there are several options.

The following are unnumbered displayed equations representing the Fourier Transform and its inverse using the modern \texttt{\textbackslash[...\textbackslash]} syntax.
\[
\mathcal{F}\{f(t)\}=F(\omega)=\int_{-\infty}^{\infty} f(t)\,e^{-i\omega t}\,dt
\]
\[
f(t)=\frac{1}{2\pi}\int_{-\infty}^{\infty} F(\omega)\,e^{i\omega t}\,d\omega
\]

For a single numbered equation, use the \texttt{equation} environment. You can add a \texttt{\textbackslash label} to reference it later in your text.
\begin{equation}
    \label{eq:euler}
    e^{i\pi} + 1 = 0
\end{equation}

For multiple aligned equations, the \texttt{align} environment. Use \texttt{\&} to mark the alignment point.
\begin{align}
    \label{eq:maxwell1}
    \nabla \cdot \mathbf{E} &= \frac{\rho}{\varepsilon_0} \\
    \label{eq:maxwell2}
    \nabla \cdot \mathbf{B} &= 0 \\
    \label{eq:maxwell3}
    \nabla \times \mathbf{E} &= -\frac{\partial \mathbf{B}}{\partial t} \\
    \label{eq:maxwell4}
    \nabla \times \mathbf{B} &= \mu_0 \left( \mathbf{J} + \varepsilon_0 \frac{\partial \mathbf{E}}{\partial t} \right)
\end{align}
You can refer to Euler's identity as Equation~\ref{eq:euler} and the first Maxwell equation as Equation~\ref{eq:maxwell1}.
\subsection{Source code}
You can include source code snippets using the \texttt{code} environment. Specify the programming language as an argument to enable syntax highlighting. For example, here is a Python code snippet:
\begin{code}{python}
print("Hello world!")
for i in range(3):
    print(i)
\end{code}

Here is a C code snippet demonstrating a simple ``Hello, world!'' program along with a function to test long line wrapping:
\begin{code}{c}
  #include <stdio.h>

  int main(void) {
      printf("Hello, world!\n");
    return 0;
  }
  void long_line_test(void){printf("LONG_LINE_TEST: 0123456789ABCDEFGHIJKLMNOPQRSTUVWXYZabcdefghijklmnopqrstuvwxyz--REPEAT--0123456789ABCDEFGHIJKLMNOPQRSTUVWXYZabcdefghijklmnopqrstuvwxyz :: punctuation !@#$%^&*()_+-={}[]|:;\"'<>,.?/ \\ \\ \\ testing very very very very very very long single line for wrap behavior in LaTeX code environment end.\n");}
\end{code}